\section{Zadanie 5 (Kacper Orszulak)}
Niech $X, Y$ będą niezależnymi zmiennymi o takich samych rozkładach i skończonej wartości oczekiwanej oraz wariancji. Niech $U = X+Y$ i $V = X-Y$.

\begin{enumerate}[label=(\roman*), itemsep=0pt, topsep=0pt, partopsep=0pt, parsep=0pt]
    \item Wykaż, że $U, V$ są nieskorelowane.
    \item Podaj przykład $X, Y$ takich, że $U, V$ nie są niezależne.
    \item Udowodnij, że jeśli $X, Y$ mają rozkład $N(0, 1)$, to $U, V$ są niezależne.
\end{enumerate}

\subsection*{Rozwiązanie}
\subsubsection*{(i)}
Z definicji zmienne $U, V$ są skorelowane wtw. $Cov(U, V) = 0$. Bardziej użyteczny będzie natomiast równoważny warunek: $\mathbb{E}[UV] = \mathbb{E}[U] \mathbb{E}[V]$. Sprawdźmy go zatem:
\[
    \mathbb{E}[UV] = \mathbb{E}[(X+Y)(X-Y)] = \mathbb{E}[X^2 - Y^2] = \mathbb{E}[X^2] - \mathbb{E}[Y^2]
\]
\[
    \mathbb{E}[U] \mathbb{E}[V] = \mathbb{E}[X+Y] \mathbb{E}[X-Y] = (\mathbb{E}[X] + \mathbb{E}[Y])(\mathbb{E}[X] - \mathbb{E}[Y]) = \mathbb{E}[X]^2 - \mathbb{E}[Y]^2
\]
Ostatnie porównanie:
\begin{align*}
    \mathbb{E}[X^2] - \mathbb{E}[Y^2] &= \mathbb{E}[X]^2 - \mathbb{E}[Y]^2 \\
    \mathbb{E}[X^2] - \mathbb{E}[X]^2 &= \mathbb{E}[Y^2] - \mathbb{E}[Y]^2 \
\end{align*}
Oczywiście jest to prawda, ponieważ obie zmienne mają ten sam rozkład. $\hfill \square$

\subsubsection*{(ii)}
Niech $X, Y$ to niezależne zmienne losowe o tym samym rozkładzie t. że
\[
    \mathcal{P}(X=-1)=\mathcal{P}(X=+1)=\frac{1}{2}
\]
Wtedy oczywiście
\begin{align*}
    \mathcal{P}(X+Y=-2) = \mathcal{P}(X-Y=-2) &= \frac{1}{4} \\
    \mathcal{P}(X+Y=-2 \; \land \; X-Y=-2) = 0 &\not = \frac{1}{4}
\end{align*}
Znaleźliśmy więc szukany kontrprzykład.

\subsubsection*{(iii)}
Niech $X,Y \sim N(0, 1)$ będą niezależne. Dowodzimy, że $U, V$ są niezależne.

Oczywiście z własności rozkładu normalnego $-Y \sim N(0, 1)$. Dla przypomnienia $\substack{U := X+Y \\ V := X - Y}$. Z twierdzenia z wykładu wiemy zatem, że $U, V \sim N(0, 2)$.

\textbf{Szukamy równości: } $\mathcal{P}(U \le u_0 \land V \le v_0) = \mathcal{P}(U \le u_0)\mathcal{P}(V \le v_0)$

\underline{Prawa strona równości}, to po prostu:
\[
    \mathcal{P}(U \le u_0 \land V \le v_0) = \left(\int_{-\infty}^{u_0} f_U(u) du\right)\left(\int_{-\infty}^{v_0} f_V(v) dv\right).
\]

Rozpisujemy \underline{lewą stronę równości}. Aby to zrobić posłużymy się zamianą współrzędnych znaną z analizy. Przypomnijmy ogólną formułę:
\[
    \int_A f = |\Phi'| \int_{\Phi^{-1}(A)} (f \circ \Phi)
\]
Jak to rozumieć? Pomysł jest taki, że chcemy policzyć całkę funkcji $f$ na jakimś ``brzydkim'' zbiorze $A$. Okazuje się jednak, że znamy przekształcenie $\Phi$ (intuicyjnie wystarczy by to było coś podobnego do bijekcji - w szczególności bijekcja jest dobra), które zamienia pewien prostokąt (po którym oczywiście łatwo całkować) na tę ``brzydką'' przestrzeń. W tym zapisie $\Phi^{-1}(A)$ oznacza zatem przeciwobraz ``brzydkiego'' zbioru, a więc nasz docelowy prostokąt.

Niech
\[
    \Phi: \; \mathbb{R}^2 \ni (u, v) \to (x, y) = \left(\frac{u+v}{2}, \frac{u-v}{2}\right) \in \mathbb{R}^2.
\]

Jacobian takiego przekształcenia liniowego to

\[
    J_{\Phi} = \Phi' =
    \begin{bmatrix}
        \frac{\partial x}{\partial u} &
        \frac{\partial x}{\partial v} \\[4pt]
        \frac{\partial y}{\partial u} &
        \frac{\partial y}{\partial v} \\[4pt]
    \end{bmatrix} =
    \begin{bmatrix}
        \frac{1}{2} &
        \frac{1}{2} \\[4pt]
        \frac{1}{2} &
        -\frac{1}{2} \\[4pt]
    \end{bmatrix}
\]
\[
    \det \Phi' = -\frac{1}{4} -\frac{1}{4} = -\frac{1}{2}
\]
\[
    |\Phi'| = |\det \Phi'| = \frac{1}{2}.
\]

Definiujemy $A$ jako taki zbiór punktów $(x, y)$, które spełniają nierówność lewej strony:
\[
    A := \left\lbrace (x, y): \; \begin{cases} x+y \le u_0 \\ x-y \le v_0 \end{cases} \right\rbrace
\]
Oczywiście takiemu zbiorowi odpowiada prostokąt
\[
    \Phi^{-1}(A) = \left\lbrace (u, v): \; \begin{cases} u \le u_0 \\ v \le v_0 \end{cases} \right\rbrace.
\]

Możemy wreszcie przejść do właściwego rozpisywania lewej strony równości:
\begin{align*}
    \mathcal{P}(U \le u_0 \land V \le v_0) &= \mathcal{P}((x, y) \in A) \\
    &= \iint_A f_{X,Y}(x, y) dy \; dx \\
    &= |\Phi'| \iint_{\Phi^{-1}(A)} (f_{X,Y} \circ \Phi)(u, v) \; dv \; du \\
    &= \frac{1}{2} \int_{-\infty}^{u_0} \int_{-\infty}^{v_0} f_{X,Y}\left(\frac{u+v}{2}, \frac{u-v}{2}\right) \; dv \; du
\end{align*}

Skorzystamy teraz z niezależności $X, Y$ oraz z następujących wzorów:
\begin{align*}
    \phi(x) &= \frac{1}{\sqrt{2\pi}\sigma} exp\left(-\frac{1}{2}\left(\frac{x-\mu}{\sigma}\right)^2\right) \\
    f_X(x) = f_Y(x) = \phi(x) &= \frac{1}{\sqrt{2\pi}} exp\left(-\frac{x^2}{2}\right) \\
    f_U(u) = f_V(u) = \phi\left(\frac{u}{\sqrt{2}}\right) &= \frac{1}{2\sqrt{\pi}} exp\left(-\frac{u^2}{4}\right)
\end{align*}

Rozpisujemy:
\begin{align*}
    f_{X,Y}\left(\frac{u+v}{2}, \frac{u-v}{2}\right)
    &= f_{X}\left(\frac{u+v}{2}\right) f_{Y}\left(\frac{u-v}{2}\right) \\
    &= \frac{1}{\sqrt{2\pi}} exp\left(-\frac{1}{2} \left(\frac{u+v}{2}\right)^2\right) \frac{1}{\sqrt{2\pi}} exp\left(-\frac{1}{2} \left(\frac{u-v}{2}\right)^2\right) \\
    &= \frac{1}{2\pi} exp\left( -\frac{1}{2} \left( \left( \frac{u+v}{2}\right)^2 + \left( \frac{u-v}{2}\right)^2 \right) \right) \\
    &= \frac{1}{2\pi} exp\left( - \frac{u^2 + v^2}{4} \right) \\
    &= 2 \frac{1}{2\sqrt{\pi}} \frac{1}{2\sqrt{\pi}} exp\left( - \frac{u^2}{4} \right) exp\left( - \frac{v^2}{4} \right) \\
    &= 2 f_U(u) f_V(v)
\end{align*}

Możemy teraz dopiąć naszą równość:
\begin{align*}
    \mathcal{P}(U \le u_0 \land V \le v_0) = \dots &= \frac{1}{2} \int_{-\infty}^{u_0} \int_{-\infty}^{v_0} f_{X,Y}\left(\frac{u+v}{2}, \frac{u-v}{2}\right) \; dv \; du \\
    &= \frac{1}{2} \int_{-\infty}^{u_0} \int_{-\infty}^{v_0} 2 f_U(u) f_V(v) \; dv \; du \\
    &= \int_{-\infty}^{u_0} f_U(u) \left( \int_{-\infty}^{v_0} f_V(v) \; dv \right) du \\
    &= \left( \int_{-\infty}^{u_0} f_U(u) \; du \right) \; \left( \int_{-\infty}^{v_0} f_V(v) \; dv \right) \\
    &= \mathcal{P}(U \le u_0 \land V \le v_0)
\end{align*}
$\hfill \square$
